\uuid{K8Xj}
\titre{Chasles/changements de variables}
\auteur{hanani ; mzari}
\organisation{exo7}
\licence{CC BY-NC-SA 4.0}
\datecreate{2020-01-20}
\niveau{L1}
\module{Analyse}
\chapitre{Calcul d'intégrales}
\sousChapitre{Autre}
\difficulte{2}
\theme{Calculs d'intégrales, Chasles, changements de variables}
\exo7qcmid{lille-2.integrales.011}

\contenu{
\begin{question}
\qclassification{127.99}
\qtags{motcle=Chasles/changements de variables}
L'intégrale $\displaystyle \int _{-\pi/6}^{\pi/3}\tan x\, \mathrm{d}x$ est égale à :
\begin{answers}
\bad{$\displaystyle \frac{4}{3}$.}
\bad{$\displaystyle \frac{2\sqrt{3}}{3}$.}
\good{$\displaystyle \int _{\pi/6}^{\pi/3}\tan x\, \mathrm{d}x$.}
\good{$\displaystyle \frac{1}{2}\ln 3$.}
\end{answers}
\vskip3mm
\begin{explanations} La relation de Chasles donne
$\displaystyle \int _{-\pi/6}^{\pi/3}\tan x\, \mathrm{d}x=\int _{-\pi/6}^{\pi/6}\tan x\, \mathrm{d}x+\int _{\pi/6}^{\pi/3}\tan x\, \mathrm{d}x$ et $\displaystyle \int _{-\pi/6}^{\pi/6}\tan x\, \mathrm{d}x=0$ car $\tan $ est impaire. Ensuite, on écrit $\displaystyle \tan x=\frac{\sin x}{\cos x}$ et on pose $t=\cos x\Rightarrow \mathrm{d}t=-\sin x\, \mathrm{d}x$. D'où,
$\displaystyle \int _{\pi/6}^{\pi/3}\tan x\, \mathrm{d}x=-\int _{\sqrt{3}/2}^{1/2}\frac{\mathrm{d}t}{t}=\ln \sqrt{3}$.
\end{explanations}
\end{question}
}
